\section{Completion of the metric classification}
\label{completion:section}

Theorem~\ref{thm:comparison} reduces metric rigidity to a volume
upper bound at the same diameter. The intrinsic quaternionic bound
has now been proved. We first check the precise classical inputs
for the complex and plane cases, then assemble the model list.
For $\HP^2$, the classical plane argument is an alternative to the
intrinsic bound, which already includes every quaternionic rank $r\ge2$.

\begin{corollary}[Complex-projective type]\label{rig:complex}
A Blaschke manifold of complex-projective type is isometric
to its equally scaled standard complex projective model.
\end{corollary}

\begin{proof}
Write its real dimension as $2r$, with $r\ge2$; the projective
line is a sphere. Its integral cohomology ring is
$\Z[u]/(u^{r+1})$, where $u$ has degree two, and it is simply
connected. Represent $u$ by a map to $\CP^\infty$
\citep[Theorem~4.57]{Hatcher2002}. A closed
smooth $2r$-manifold has the homotopy type of a finite
$2r$-dimensional CW complex, so cellular approximation
\citep[Theorem~4.8]{Hatcher2002} places the map in the
$2r$-skeleton $\CP^r$. Pullback identifies the integral
cohomology generators and all their powers. The universal
coefficient theorem for these finite CW complexes then gives
an integral homology isomorphism. The homology Whitehead
theorem for simply connected CW complexes
\citep[Corollary~4.33]{Hatcher2002} makes this map a homotopy
equivalence. Yang's volume theorem
\citep[Theorem, p.~380]{Yang1991} therefore gives equality with the model at the
same prime geodesic period, equivalently the same diameter.
Corollary~\ref{thm:volume-reduction} applies.
\end{proof}

\begin{corollary}[Quaternionic and Cayley planes]\label{rig:planes}
A Blaschke manifold with model $\HP^2$ or $\OP^2$ is isometric
to that model at the same diameter.
\end{corollary}

\begin{proof}
The projective-plane structure of such a Blaschke manifold is
given by \citet[Lemma~5]{McKay2016}. The classification of smooth
topological projective planes
\citep[Corollary~C]{KramerStolz2007} makes $M$ diffeomorphic to
its standard model. Both metrics are Zoll: all their unit-speed
geodesics are simply closed with the same prime length $T=2D$.
After a common rescaling to $T=\pi$, the Riemannian weak Blaschke
volume theorem recorded by Wilking and attributed there to
Reznikov applies: a manifold with this common least period
that is homeomorphic to a compact rank-one symmetric space has
the same volume as its model \citep[p.~282]{Wilking2001}.
The diffeomorphism above supplies its homeomorphism hypothesis.
Thus $V=V_0$, and this equality persists when the common rescaling
is undone. Corollary~\ref{thm:volume-reduction} applies.
\end{proof}

\begin{corollary}[Quaternionic type]\label{rig:quaternionic}
A Blaschke manifold of quaternionic type is isometric to its
standard quaternionic model at the same diameter.
\end{corollary}

\begin{proof}
For $\HP^r$ with $r\ge2$, use Theorem~\ref{alg:upper} and
Proposition~\ref{geo:contact}. At the common period $T=\pi$, the
contact formula has the same positive volume factor for $M$ and its
model, so
\[
 \frac{N}{N_0}=\frac{\Vol(M,g)}{\Vol(\HP^r)}.
\]
Here $N_0$ is the model number in \eqref{geo:model-number}, and the
model volume computed there is
\[
 \Vol(\HP^r)=\frac{\pi^{2r}}{(2r+1)!}.
\]
Theorem~\ref{alg:upper} gives $N\le N_0$, hence
$\Vol(M,g)\le\Vol(\HP^r)$.
Corollary~\ref{thm:volume-reduction} proves the metric statement.
The projective line is already the sphere case.
\end{proof}

\begin{proof}[Proof of Theorem~\ref{thm:blaschke}]
Lemma~\ref{geom:cut-time} gives the common minimizing distance
from the stated global injectivity-radius equality. The
one-dimensional, sphere-type and real-projective-type cases
are covered by Section~\ref{geom:section}. The classical
structure list leaves complex type, quaternionic type, and the
Cayley plane. These are respectively
Corollaries~\ref{rig:complex}, \ref{rig:quaternionic},
and~\ref{rig:planes}. Each conclusion is a metric isometry
with the model at equal diameter. Undoing the normalization
therefore gives exactly the asserted classification up to scale.
\end{proof}
